\chapter{Magnetic entropy and magnetocaloric effect (menu 41)}
\label{ch:menu41}
\menulabel{41}

\section{Purpose}

The magnetic entropy $S(T)$ and the isothermal entropy change
$\Delta S(T, H)$ are the figures of merit of the magnetocaloric effect.
This menu computes them from ORCA data by two routes, chosen by content:
the spin-orbit level spectrum (partition function) or the POLY\_ANISO
magnetization table (Maxwell relation).

\section{What you need}

Either a SINGLE\_ANISO output (or POLY\_ANISO per-center spectra) with the
``Spin-orbit energy spectra'' block, or a POLY\_ANISO run with
\code{HINT}/\code{TMAG} (the ``HIGH-FIELD POWDER MAGNETIZATION'' table).

\section{Walkthrough}

The example reads the POLY\_ANISO M(H) probe (6 temperatures $\times$ 71
fields, 0--7~T) and reports the $-\Delta S$ table with the probe maximum
-- 10.8404~J~mol$^{-1}$~K$^{-1}$ at $T$ = 1.90~K and 7~T -- then reads the
menu-36 Co spectrum through the levels route (a degenerate ground doublet
gives the flat $S = R\ln 2$ branch; the truncated level list caps the
useful range).

\lstinputlisting[style=fbkcode, firstline=53, lastline=69,
  title={The menu-41 example (the Maxwell route on the magnetization
  table, with the probe maximum)}]
  {examples/expected/m41_magnetocaloric.txt}

\section{What you get}

The report \code{<output>.mce.fbk.md}: levels route -- $S_{\mathrm{mag}}(T)$ at
the summary temperatures plus the $R\ln N$ high-temperature check (an
upper bound, since the level list is truncated); magnetization route --
the $-\Delta S(T, H)$ table at probe fields with the maximum, from the
Maxwell relation.

Beside the report, the menu writes a plot-ready companion
\code{<output>.mce.fbk.csv}, one shape per route: the levels route's
$S_{\mathrm{mag}}(T)$ grid
(\code{temperature\_K,entropy\_J\_per\_mol\_per\_K}), or the Maxwell route's
$-\Delta S(T, H)$ table in tidy long form over the full printed field grid
(\code{T\_mid\_K,field\_T,minus\_delta\_S\_J\_per\_mol\_per\_K}) -- the
report shows five probe fields, the companion carries them all
(Chapter~\ref{ch:formats}).

\section{Conventions, cross-checks and boundaries (measured)}

\begin{readbox}
\begin{itemize}
  \item \textbf{Conventions follow Szalowski \& Kowalewska 2020}: the
    Gibbs-function route $S = R(\ln Z + \langle E\rangle/kT)$ (their
    Eq.~(4)--(5)) and $\Delta S_T = S(T, B_i) - S(T, B_f)$ (their
    Eq.~(11)); the report prints $-\Delta S$ (positive = direct MCE).
  \item \textbf{The levels route is $B = 0$} (the printed levels are
    field-free) and each printed line counts as one state (the engine's
    level list handles degeneracies; e.g.\ the Co doublet gives
    $R\ln 2$).
  \item \textbf{The truncated level list bounds the range}: the
    high-temperature check $R\ln N$ is an upper bound -- do not trust $S$
    above the saturation temperature (measured: the Co spectrum's
    $R\ln 4$ extreme is never approached below 300~K because the excited
    doublet lies at 29361~cm$^{-1}$).
  \item \textbf{The Maxwell route needs converged data} (a fine
    temperature grid); the trapezoid anchors at $H = 0$ with
    $(\mathrm{d}M/\mathrm{d}T) = 0$; the engine's field grid starts at
    0.0001~T.
  \item \textbf{Probe boundary}: the shipped magnetization fixture is the
    two-center flow probe (invented $J$), so its $-\Delta S$ values are
    not a physical result -- the format and the method chain are what it
    demonstrates.
\end{itemize}
\end{readbox}
